% Standalone supplementary material for the QCE26 NIER short paper.
% Named build:     latexmk -pdf -interaction=nonstopmode -halt-on-error supplement.tex
% Anonymous build: latexmk -pdf -interaction=nonstopmode -halt-on-error supplement_anon.tex
\documentclass[10pt,conference]{IEEEtran}

\usepackage[T1]{fontenc}
\usepackage[utf8]{inputenc}
\usepackage{microtype}
\usepackage{amsmath,amssymb}
\usepackage{graphicx}
\graphicspath{{gfx/}{figures/}}
\usepackage{booktabs}
\usepackage{array}
\usepackage{balance}
\usepackage{cite}
\usepackage{xcolor}
\usepackage[hypertexnames=false]{hyperref}
\definecolor{paperblue}{HTML}{0057FF}
\definecolor{citegreen}{HTML}{39FF14}
\hypersetup{%
  colorlinks=true,
  linkcolor=paperblue,
  citecolor=citegreen,
  urlcolor=paperblue,
  pdfborder={0 0 0},
  pdftitle={Supplementary Material: From Raw Shots to Expectation Values: Shot Grouping Improves Time Series Prediction},
  pdfsubject={Robust protocol, exact covariance identity, ablations, hardware provenance, and reproducibility}
}
\urlstyle{same}

\newcommand{\nrmse}{\mathrm{NRMSE}}
\newcommand{\EV}{\mathrm{EV}}
\newcommand{\Grouped}{\mathrm{Grouped}}
% Manuscript-facing outcomes from the evidence-preserving merge in
% results/qai_robust_protocol_santa_fe_manuscript/.
\newcommand{\QAIFeatureDimension}{11}
\newcommand{\QAIFeatureEffectiveRank}{5.15}

\newcommand{\QAIMackeyEVGrouped}{0.0445}
\newcommand{\QAIMackeyEVGroupedLow}{0.0372}
\newcommand{\QAIMackeyEVGroupedHigh}{0.0519}
\newcommand{\QAISantaFeEVGrouped}{0.0848}
\newcommand{\QAISantaFeEVGroupedLow}{0.0174}
\newcommand{\QAISantaFeEVGroupedHigh}{0.1522}
\newcommand{\QAIRosslerEVGrouped}{0.0349}
\newcommand{\QAIRosslerEVGroupedLow}{0.0139}
\newcommand{\QAIRosslerEVGroupedHigh}{0.0559}
\newcommand{\QAISunspotsEVGrouped}{0.0360}
\newcommand{\QAISunspotsEVGroupedLow}{0.0103}
\newcommand{\QAISunspotsEVGroupedHigh}{0.0616}

% Protocol-matched Santa Fe hardware extension. The intervals are descriptive
% paired intervals over twenty disjoint shot blocks within one acquisition.
\newcommand{\QAIHardwareFezSantaEV}{0.750}
\newcommand{\QAIHardwareFezSantaGrouped}{0.711}
\newcommand{\QAIHardwareFezSantaGain}{0.038}
\newcommand{\QAIHardwareFezSantaGainLow}{0.027}
\newcommand{\QAIHardwareFezSantaGainHigh}{0.050}
\newcommand{\QAIHardwareMarrakeshSantaEV}{0.744}
\newcommand{\QAIHardwareMarrakeshSantaGrouped}{0.713}
\newcommand{\QAIHardwareMarrakeshSantaGain}{0.031}
\newcommand{\QAIHardwareMarrakeshSantaGainLow}{0.023}
\newcommand{\QAIHardwareMarrakeshSantaGainHigh}{0.039}
\newcommand{\QAIHardwareSantaWins}{19/20}
\newcommand{\QAIHardwareFezSeconds}{55}
\newcommand{\QAIHardwareMarrakeshSeconds}{56}
\newcommand{\QAIHardwareTotalSeconds}{111}

\setlength{\textfloatsep}{8pt plus 1pt minus 2pt}
\setlength{\dbltextfloatsep}{8pt plus 1pt minus 2pt}
\setlength{\floatsep}{7pt plus 1pt minus 1pt}

\begin{document}

\title{Supplementary Material for\\
\emph{From Raw Shots to Expectation Values:\\
Shot Grouping Improves Time Series Prediction}}

\ifdefined\ANON
  \hypersetup{pdfauthor={}}
  \author{\IEEEauthorblockN{Anonymous Authors}}
\else
  \hypersetup{pdfauthor={Markus Baumann; Maximilian Zorn; Thomas Gabor; Claudia Linnhoff-Popien; Jonas Stein}}
  \author{%
  \IEEEauthorblockN{Markus Baumann, Maximilian Zorn, Thomas Gabor,
  Claudia Linnhoff-Popien, and Jonas Stein}}
\fi

\maketitle
\bstctlcite{BSTcontrol}

\paragraph{Scope and artifact status.}
This document records the exact algebra, shared four-task protocol, completed
simulation, and provenance behind the main paper. The manuscript evidence
contains 13,180 replicate rows and 9,720 aggregation-spectrum rows for
Mackey--Glass, Santa Fe laser, R\"ossler, and sunspots. All four simulation
tasks use the same six-qubit reservoir, features, budgets, chronological joint
tuning, parameter grids, seed design, and test-access rule; separate source
manifests preserve their data and implementation provenance. At the balanced
setting, Grouped is supported on
Santa Fe, R\"ossler, and sunspots, and positive but inconclusive on
Mackey--Glass. The completed unmitigated hardware transfer applies a separate
fixed ten-qubit protocol to the same four tasks on two processors. Its block-level
intervals describe shot variation within one acquisition, not independent
hardware campaigns.

\appendices

\section{Exact Identity and Implementation Scaling}
\label{app:identity}

For aggregation scale $k$, let $G=N/k$ and let $x_t^{(g)}\in\mathbb{R}^d$ denote the
leaked and lagged feature vector of group $g$ at training time $t$.  With an
unpenalized intercept $b$, every time step receives total weight one:
\begin{equation}
  \mathcal L_k(b,w)=\frac{1}{|\mathcal T|G}
  \sum_{t\in\mathcal T}\sum_{g=1}^{G}
  \bigl(y_t-b-w^\top x_t^{(g)}\bigr)^2
  +\lambda\lVert w\rVert_2^2.
  \label{eq:supp-grouped}
\end{equation}
Define $\bar x_t=G^{-1}\sum_gx_t^{(g)}$ and
\begin{equation}
  C_k=\frac{1}{|\mathcal T|G}\sum_{t,g}
  (x_t^{(g)}-\bar x_t)(x_t^{(g)}-\bar x_t)^\top .
  \label{eq:supp-cov}
\end{equation}

\paragraph{Proof of the exact decomposition.}
Write $x_t^{(g)}=\bar x_t+d_t^{(g)}$, where
$\sum_g d_t^{(g)}=0$.  Expanding the square in
\eqref{eq:supp-grouped}, the cross term is proportional to
$(y_t-b-w^\top\bar x_t)w^\top\sum_g d_t^{(g)}$ and therefore vanishes at every
$t$; the remaining deviation term is
$G^{-1}\sum_g(w^\top d_t^{(g)})^2=w^\top C_{k,t}w$.  Averaging over $t$ gives
\begin{equation}
  \mathcal L_k(b,w)=\frac{1}{|\mathcal T|}\sum_t
  (y_t-b-w^\top\bar x_t)^2+w^\top C_kw
  +\lambda\lVert w\rVert_2^2,
  \label{eq:supp-identity}
\end{equation}
which is deterministic and assumes neither Gaussian noise, isotropy, nor
independence across time.

\subsection{Scaling used in the implementation}

For every fold and candidate $k$, standardization is fit on grouped training
rows only.  The same transformed rows define both $C_k$ and the grouped fit.
The implementation uses \texttt{sklearn.linear\_model.Ridge}, whose loss is a
sum rather than the mean loss in \eqref{eq:supp-grouped}.  It therefore passes
\begin{equation}
  \texttt{alpha}_{\rm sklearn}=\lambda\,|\mathcal T|G.
  \label{eq:supp-sklearn-scaling}
\end{equation}
This preserves the meaning of $\lambda$ as $G$ changes.  The generalized-ridge
solver independently reconstructs the solution from $\bar x_t$, $C_k$, and
$\lambda I$; unit tests require its predictions to match direct grouped Ridge
to numerical precision.

The directly tuned Cov-EV baseline is intentionally different.  It fits
EV rows with
$\alpha w^\top\widehat\Sigma_{\rm shot}w+\lambda\lVert w\rVert_2^2$ and
tunes $(\alpha,\lambda)$ on the same chronological folds. It does not inherit
Grouped training's $k$, $C_k$, or tuned penalty and is therefore a genuine competing
baseline rather than an identity check.  The decomposition is consistent with
the general connection between feature noise and Tikhonov regularization
\cite{bishop1995training}, but here it holds for the observed finite record
exactly.

\section{Frozen Evaluation Protocol}
\label{app:protocol}

The primary experiment was frozen in
\path{configs/qai_robust_protocol.json} before inspecting redesigned forecast
outcomes.  All methods reuse the same sampled local-$Z$ measurement record.  The classical
baseline receives the scalar input with the identical leak $\eta=0.2$ and ten
causal lags.  The primary quantum representation contains six local $Z_i$ and
five neighboring $Z_iZ_{i+1}$ features reconstructed from the same shots.

\begin{table}[t]
\caption{Three predeclared settings.  Every record uses exactly $TN=10{,}000$
circuit executions.}
\label{tab:supp-budget}
\centering
\begin{tabular}{lrrl}
\toprule
Setting & $T$ & $N$ & candidate $k$ \\
\midrule
Shot rich & 250 & 40 & 1,2,4,5,8,10,20,40 \\
Balanced  & 500 & 20 & 1,2,4,5,10,20 \\
Data rich & 1000& 10 & 1,2,5,10 \\
\bottomrule
\end{tabular}
\end{table}

\subsection{Chronological boundaries and normalization}

Indices in Table~\ref{tab:supp-folds} are zero-based, half-open intervals.
After washout $t_0=30$, the last 20\% of the usable sequence is the untouched
test block.  The first half of the pre-test interval initializes an expanding
training prefix; the remainder is split into three consecutive blocked
development folds. Input normalization is frozen using raw series values only
through the first training stop (inclusive because the next-step target is
formed from the following value).  One train-fitted affine min--max map is used
on both branches.  Circuit inputs alone are clipped to $[0,1]$ for angle
encoding; next-step targets retain the same affine scale without clipping and
may therefore lie below zero or above one.  Feature standardizers are then fit
separately for every fold and candidate using training rows only.

\begin{table*}[t]
\caption{Exact fold boundaries generated by the frozen protocol.  Each entry
is $[\mathrm{start},\mathrm{stop})$.}
\label{tab:supp-folds}
\centering
\begin{tabular}{lccccc}
\toprule
Setting & normalization stop & fold 1: fit / develop & fold 2: fit / develop & fold 3: fit / develop & final test \\
\midrule
Shot rich & 118 & [30,118) / [118,147) & [30,147) / [147,177) & [30,177) / [177,206) & [206,250) \\
Balanced  & 218 & [30,218) / [218,281) & [30,281) / [281,343) & [30,343) / [343,406) & [406,500) \\
Data rich & 418 & [30,418) / [418,547) & [30,547) / [547,677) & [30,677) / [677,806) & [806,1000) \\
\bottomrule
\end{tabular}
\end{table*}

For each candidate pair, development NRMSE is computed from one forecast per
time step after averaging the group predictions. Mean development NRMSE across
the three folds determines the tuned pair; exact ties prefer smaller $\lambda$,
then smaller $k$ (or smaller $\alpha$ for Cov-EV). The tuned configuration
is refit on $[30,t_{\rm test})$, and the final test block is evaluated once.

\subsection{Grids and nested uncertainty}

The shared ridge grid is
$\lambda\in\{10^{-3},10^{-2},10^{-1},1,10,10^2,10^3\}$.  Cov-EV uses
$\alpha\in\{0,0.01,0.03,0.1,0.3,1,3,10\}$.  Grouped searches all divisors of
$N$, including $k=1$ and $k=N$ as safe endpoint fallbacks.  The compared
methods are: matched classical input, ideal statevector QRC, EV, adaptive
Grouped, Raw, and tuned Cov-EV.

Five independent initial conditions are used for each synthetic task. The two
measured series use five windows with disjoint test blocks: sunspot origins
are separated by 240 months, and Santa Fe origins by 1,001 samples. Santa Fe
uses only the 9,093-point continuation file, so its complete input windows are
also disjoint at every budget. Reservoir seeds
$\{1729,2718,31415\}$, shot seeds $\{1009,2003,3001\}$, and primary grouping
seeds $\{17,43,97\}$ are nested factors. The inferential unit is the synthetic
trajectory or measured-series window: nested repeats are averaged within each unit,
and 95\% $t$ intervals are computed over the five unit-level means.  Paired
method differences are formed within the same unit before interval estimation;
shot or grouping seeds are never counted as independent forecasting samples.
All forecasts predict the next step from the observed input at
time $t$ is available when predicting $t+1$.  We do not evaluate autonomous
multistep rollouts.

\section{Expressive Architecture and Target-Free Diagnostics}
\label{app:architecture}

The redesigned reservoir was chosen using ideal feature behavior and hardware
depth only, before forecast targets were examined.  It uses six qubits, two
data-reuploading layers, and seed 1729.  Each layer applies input-dependent
$R_Y$ and $R_Z$ rotations to every qubit, heterogeneous nearest-neighbor
$R_{ZZ}$ rotations on a line, and fixed heterogeneous local $R_XR_YR_Z$
mixing.  This is a compact data-reuploading construction
\cite{perezsalinas2020data}; no trainable quantum parameters are used.

\begin{table}[t]
\caption{Frozen circuit identity.}
\label{tab:supp-hashes}
\centering
\begin{tabular}{p{0.28\columnwidth}p{0.61\columnwidth}}
\toprule
Field & Frozen value \\
\midrule
Circuit version & \texttt{expressive\_reuploading\_v2} \\
Qubits / layers & 6 / 2 \\
Primary seed & 1729 \\
Observables & $Z_0,\ldots,Z_5$ and $Z_iZ_{i+1}$, $i=0,\ldots,4$ \\
Parameter hash & \scriptsize\texttt{fdff90afa7f488026e7c822f33e424682}\newline
  \texttt{c2c084653d3ccd245f91e2e90038505} \\
Architecture hash & \scriptsize\texttt{d44fac7ba013f8337e293d9e08ca3b3b}\newline
  \texttt{d85be634520d7b9f1953be296770b964} \\
\bottomrule
\end{tabular}
\end{table}

The complete double-precision arrays for the two sets of six input-$Y$
weights, input-$Z$ weights, $2\times6\times3$ local angles, and $2\times5$
couplings are serialized in
\path{results/expressive_reservoir/ideal_diagnostics.json}. Hash validation is
mandatory in the simulation runner. Regenerating angles from the seed is convenient, but
the serialized values and hash, rather than a library-specific random-number
implementation, define the circuit.

\subsection{Pre-forecast validity gate}

The diagnostic evaluates 257 inputs using exact statevectors and uses no
forecast target or shot sampling.  All 11 features vary with input and all 11
leave more than 5\% variance unexplained by an affine function of the input.
The standardized feature matrix has numerical rank 11 and effective rank
5.15.  Removing the second layer's couplings changes every measured feature;
the maximum absolute effect is 1.58.  Finally, after applying the same leak and
ten-lag map, projection onto the full linear span of the filtered scalar input
leaves relative residual 0.786.  These values pass the frozen thresholds in
the diagnostic JSON.  They establish only that the measured representation is
input-dependent, nonlinear, non-collapsed, and not a linear copy of the matched
classical preprocessing.  They are not a forecast result or a quantum-advantage
claim.

\section{Completed Simulation Outcome}
\label{app:outcome}

The manuscript filter retains only the expressive $Z+ZZ$ representation and
the balanced setting. Diagnostic and ablation rows cannot enter the main
report. Nested reservoir, shot, and grouping
seeds are averaged within each of five trajectories or real-series windows before
the method means and paired Student-$t$ intervals are formed.

\begin{table*}[t]
\caption{Balanced-setting mean test NRMSE for the classical and ideal controls
and EV, Grouped, Raw, and Cov-EV.  Bold marks the best sampled quantum mean,
not the best value across the classical and ideal columns.}
\label{tab:supp-main-means}
\centering
\begin{tabular}{lcccccc}
\toprule
Task & Classical & Ideal QRC & EV & Grouped & Raw & Cov-EV \\
\midrule
Mackey--Glass & 0.100 & 0.178 & 0.468 & 0.454 & 0.617 & \textbf{0.447} \\
Santa Fe laser & 0.762 & 0.493 & 0.764 & 0.680 & 0.697 & \textbf{0.677} \\
R\"ossler & 0.096 & 0.054 & 0.414 & 0.379 & 0.580 & \textbf{0.367} \\
Sunspots & 0.387 & 0.474 & 0.553 & 0.517 & 0.558 & \textbf{0.515} \\
\bottomrule
\end{tabular}
\end{table*}

Table~\ref{tab:supp-contrasts} gives the two paired comparisons. A
positive EV-minus-Grouped value favors Grouped; a positive
Grouped-minus-Cov-EV value favors Cov-EV. Grouped has supported gains
on Santa Fe, R\"ossler, and sunspots, while Mackey--Glass is inconclusive. All
four lower bounds exceed the $-0.05$ practical margin and three exceed zero.
The numerical pattern is consistent with the shared protocol: Santa Fe has the
largest balanced gain, while Mackey--Glass remains the weakest contrast.

\begin{table*}[t]
\caption{Frozen balanced-setting paired effects over five independent
trajectories/windows.  Values are mean $[95\%\ \mathrm{CI}]$ in NRMSE.}
\label{tab:supp-contrasts}
\centering
\begin{tabular}{lcc}
\toprule
Task & EV $-$ Grouped & Grouped $-$ Cov-EV \\
\midrule
Mackey--Glass & $0.0138\ [-0.0023,\ 0.0300]$ & $0.0066\ [-0.0026,\ 0.0158]$ \\
Santa Fe laser & $0.0848\ [0.0174,\ 0.1522]$ & $0.0026\ [-0.0167,\ 0.0218]$ \\
R\"ossler & $0.0349\ [0.0139,\ 0.0559]$ & $0.0116\ [0.0021,\ 0.0212]$ \\
Sunspots & $0.0360\ [0.0103,\ 0.0616]$ & $0.0021\ [-0.0025,\ 0.0067]$ \\
\bottomrule
\end{tabular}
\end{table*}

The covariance-control gate passed its narrow condition: every upper bound in
the last column was below $0.05$.  Cov-EV nevertheless had the lower
mean on all four tasks and a supported advantage on R\"ossler. The gate means
only that its advantage did not exceed the frozen practical margin; it does not
show that Grouped is the uniquely effective regularizer.

The ideal-QRC control remains mixed. At the balanced setting, the
classical-minus-ideal differences were $-0.0775$ for Mackey--Glass, $0.2688$
for Santa Fe, $0.0420$ for R\"ossler, and $-0.0867$ for sunspots. The matched
classical input was significantly better on Mackey--Glass and sunspots; ideal
QRC was significantly better on Santa Fe; R\"ossler's interval crossed zero.
Feature validity therefore did not translate into uniform forecasting utility
or establish quantum advantage.

\section{Grouping Assignment and Robustness Contract}
\label{app:partitions}

A sampler returns an unordered set of bitstrings at each time step.  Treating
the first stored shot at time $t$ and the first stored shot at $t+1$ as the
same persistent shot index creates a classical bookkeeping link, not a physical
trajectory.  A fixed-group construction based on those positions could
therefore make grouping depend on an arbitrary storage order.

For each aggregation scale $k$ and grouping seed, the primary procedure independently permutes
the $N$ shot indices at every time step, divides the permuted record into
contiguous groups of size $k$, averages observables within each group, and only
then applies the linear leak and lag map to each group sequence.  Candidate
values of $k$ receive distinct deterministic permutations.  Because both
averaging and the leak/lag operator are linear,
\begin{equation}
  \frac{1}{G}\sum_g \mathcal P\!\left(\frac{1}{k}
  \sum_{n\in g}z_t^{(n)}\right)
  =\mathcal P\!\left(\frac{1}{N}\sum_{n=1}^{N}z_t^{(n)}\right),
  \label{eq:supp-mean-preserving}
\end{equation}
where $\mathcal P$ denotes the common leak-and-lag map.  Thus every grouping
has exactly the same mean feature sequence as EV; only the empirical
within-group covariance differs.

The headline design averages over grouping seeds 17, 43, and 97.  A balanced
setting subset repeats the primary construction for 20 additional seeds
101--120.  A fixed-group ablation uses seeds 17, 43, and 97 with one grouping
assignment shared across time.  The predeclared outputs are
\path{partition_robustness_summary.csv},
\path{fixed_lane_sensitivity.csv}, and
\path{selected_k_frequency.csv} in
\path{results/qai_robust_protocol_santa_fe_manuscript/}. They report trajectory/window-level
paired gains, the difference between fixed-group and per-time regrouping, and tuned-$k$
frequencies, respectively.  The complete campaign passes the strict reporter;
incomplete shards were not pooled.

The balanced frequency panel contains 135 nested tuned configurations per task
($5$ clusters $\times3$ reservoir seeds $\times3$ shot seeds $\times3$
grouping seeds).  They are descriptive choices, not 135 independent trials.
$k=1$ was eligible and tuning returned it once in the Santa Fe screen. Tuning
returned an interior $k$ in 96.3\% of Mackey--Glass runs, 94.1\% of Santa Fe runs, 96.3\%
of R\"ossler runs, and 97.0\% of sunspot runs. These high rates describe the
tuned configurations; they are not independent test outcomes.

\begin{table*}[t]
\caption{Grouping checks at the balanced setting.  Values are mean
$[95\%\ \mathrm{CI}]$ over five trajectory/window clusters.  Positive
EV-minus-Grouped favors Grouped; positive fixed-group-minus-per-time-regrouped
favors per-time regrouping.}
\label{tab:supp-partitions}
\centering
\begin{tabular}{lcc}
\toprule
Task & 20-grouping EV $-$ Grouped & Fixed group $-$ per-time regrouped \\
\midrule
Mackey--Glass & $-0.0064\ [-0.0436,\ 0.0309]$ & $-0.0021\ [-0.0165,\ 0.0124]$ \\
Santa Fe laser & $0.1103\ [-0.0146,\ 0.2351]$ & $0.0004\ [-0.0320,\ 0.0327]$ \\
R\"ossler & $0.0177\ [-0.0270,\ 0.0624]$ & $-0.0018\ [-0.0134,\ 0.0098]$ \\
Sunspots & $0.0493\ [0.0219,\ 0.0768]$ & $-0.0017\ [-0.0079,\ 0.0045]$ \\
\bottomrule
\end{tabular}
\end{table*}

Only sunspots had a positive clustered interval in the 20-grouping subset.
Every fixed-group interval crossed zero.  The result therefore supports neither
a universal grouping-robust benefit nor a meaningful advantage of one
grouping convention over the other.

\section{Completed Four-Task Hardware Transfer}
\label{app:hardware}

Mackey--Glass, Santa Fe, R\"ossler, and sunspots use the same unmitigated
ten-qubit circuit, local-$Z$ readout, leak and ten-lag postprocessing,
chronological development/test split, and tuning grids on \texttt{ibm\_fez} and
\texttt{ibm\_marrakesh}. Each task--processor pair contributes one $T=500$
acquisition with 400 ordered shots per step. The record is partitioned into
twenty disjoint $N=20$ blocks, each retaining the 10,000-execution budget used
by the balanced simulation setting.

This hardware circuit is separate from the six-qubit expressive reservoir used
in simulation. It uses one amplitude-encoded input qubit, nine
Hadamard-initialized memory qubits, a depth-one ring of
$R_{ZZ}(\pi/2)$ couplings, $R_X(\pi/4)$ rotations, and ten local-$Z$
measurements. The hardware study tests the aggregation procedure on a second
fixed reservoir realization. It is not a reproduction of the six-qubit feature
map and does not compare qubit counts. The circuit and physical layout are
fixed within each processor, so EV, Grouped, and Raw differ only in their
classical use of the same ordered measurement record.

Santa Fe uses continuation window~5, $[4587,5088)$, disjoint from the five
simulation windows. It was acquired on July~12, 2026, under the same circuit,
archived physical layouts, and analysis as the earlier hardware tasks. The Fez
job \texttt{d99mf4t2su3c739kl9k0} used \QAIHardwareFezSeconds{} QPU seconds;
the Marrakesh job \texttt{d99mf52f47jc73a99fj0} used
\QAIHardwareMarrakeshSeconds{}, for \QAIHardwareTotalSeconds{} seconds total.
On Fez, Santa EV and Grouped NRMSE are \QAIHardwareFezSantaEV{} and
\QAIHardwareFezSantaGrouped{}; on Marrakesh they are
\QAIHardwareMarrakeshSantaEV{} and \QAIHardwareMarrakeshSantaGrouped{}.
Grouped wins \QAIHardwareSantaWins{} blocks on each processor.

The paired intervals use the twenty disjoint shot blocks within each single
acquisition. They are descriptive checks of shot-record organization and must
not be interpreted as replication across calibration periods, jobs, or devices.

\section{Reproducibility Commands and Artifact Inventory}
\label{app:reproducibility}

All commands below run from the repository root.  First regenerate the
target-free diagnostic and check its serialized parameter hash:
\begin{scriptsize}
\begin{verbatim}
D=scripts/\
run_expressive_reservoir_diagnostics.py
python "$D" --feature-grid-csv \
  results/expressive_reservoir/\
ideal_feature_grid.csv
\end{verbatim}
\end{scriptsize}
Inspect the current four task job plan and verify the included evidence:
\begin{scriptsize}
\begin{verbatim}
python scripts/run_robust_evaluation.py \
  --config configs/benchmark.json \
  --results-dir results/benchmark_dry_run \
  --dry-run
python scripts/verify_santa_fe_manuscript.py
python scripts/plot_manuscript_figures.py
\end{verbatim}
\end{scriptsize}
The completed IBM acquisitions cannot be reproduced without new QPU use. The
released raw records can be reanalyzed and the hardware figure redrawn offline:
\begin{scriptsize}
\begin{verbatim}
python scripts/analyze_qai_hardware.py
python scripts/analyze_qai_hardware.py \
  --subdir marrakesh
python scripts/plot_hardware_figure.py
\end{verbatim}
\end{scriptsize}
The extension manifest preserves live availability, frozen input hashes,
physical layouts, job identifiers, ordered-bitstring hashes, and charged time;
credentials are neither stored nor released.

\subsection{Artifact inventory}

\noindent\textbf{Frozen inputs:}
The current runnable protocol is \path{configs/benchmark.json}. The files
\path{configs/qai_robust_protocol.json},
\path{configs/qai_robust_protocol_santa_fe.json},
and \path{results/expressive_reservoir/ideal_diagnostics.json} are retained
because their hashes are recorded by the completed source manifests.
\textbf{Simulation provenance:}
the source manifests in \path{results/qai_robust_protocol_full/} and
\path{results/qai_santa_fe_laser_screen/} record configuration,
implementation, commit, data, and test-access hashes. The merge manifest in
\path{results/qai_robust_protocol_santa_fe_manuscript/} records both source
hashes and output provenance. Source \path{shards/} contain checkpoints. The
merged \path{replicate_results.csv},
\path{cluster_summary.csv}, \path{cluster_paired_stats.csv}, and
\path{partition_robustness_summary.csv} contain completed nested rows, cluster
summaries, EV-paired differences, and grouping sensitivity.
The paper directory contains the balanced table source and regenerated vector
figures; the merged result directory contains the manuscript-facing tables and
their hashes.
\textbf{Hardware provenance:}
\path{results/qai_hardware_santa_fe_extension/manifest.json} records the
completed Santa acquisition and promotion into the two four-task archives under
\path{results/qai_hardware/}. The raw shots, task inputs, layouts, operational
manifests, and derived tables needed to audit the main-paper hardware figure
are released together.

Build this standalone supplement with
\begin{scriptsize}
\begin{verbatim}
cd paper
latexmk -pdf -interaction=nonstopmode \
  -halt-on-error supplement.tex
\end{verbatim}
\end{scriptsize}
The bibliography and IEEE class are the same files used by the six-page main
paper.  The named artifact is available at
\ifdefined\ANON
an anonymous review link supplied separately with the submission.
\else
\href{https://github.com/eybmits/split-ensemble-qrc}{github.com/eybmits/split-ensemble-qrc}.
\fi

\bibliographystyle{IEEEtranDoi}
\bibliography{bstcontrol,references}

\end{document}
